% =====================================================================
%  Model 2 -- how the per-fish latent works.
%  Drop-in replacement for the "Estimation" subsection.
%  Requires: amsmath, booktabs.  (No siunitx, no algorithm packages.)
%  Assumes the macros \dth and \E are already defined.
% =====================================================================

\subsubsection{How the mechanism works}

It is worth setting out the mechanism concretely, because the way the
between-fish spread is produced --- and the reason it needs a correction at all
--- is not obvious from the model equations alone.

\paragraph{What is being fitted.} Model~2 adds three groups of parameters to
Model~1:

\begin{center}
\begin{tabular}{lrl}
\toprule
 & size & role \\
\midrule
network weights $\theta$          & ${\sim}20{,}000$ & shared by all fish, exactly as Model~1 \\
latent table $Z$                  & $15\times 2$     & one row per \emph{training} fish \\
loadings $B_\theta,B_\ell,B_\tau$ & $3\times 2$      & how a latent becomes a nudge \\
\bottomrule
\end{tabular}
\end{center}

\noindent Thirty-six new numbers in total. The network itself is unchanged: no
layer takes the fish index as an input, and the shared body is evaluated
identically for every fish.

\paragraph{One training row.} For a bout belonging to fish $i$, with features
$x$ and observations $(\dth,\ell,\tau)$:
\begin{align*}
h &= \mathrm{ReLU}(Wx+b) && \text{shared body, identical for every fish}\\
z &= Z[i] && \text{look up fish $i$'s two numbers}\\
o^{\text{turn}}_F &= \text{head}_\theta(h)_0 + B_\theta z\\
\mu^\ell &= \text{head}_\ell(h,|\dth|) + B_\ell z\\
\mu^\tau &= \text{head}_\tau(h,|\dth|,\log\ell) + B_\tau z\\
\mathcal{L} &= -\log p(\dth,\ell,\tau\mid\cdot) + \tfrac{1}{2}\lVert z\rVert^2/N .
\end{align*}
Two numbers are dotted into three scalars. That is the entire forward
mechanism; the fish index enters nowhere else.

\paragraph{The two gradients.} The mechanism is best understood through what
the two parameter groups learn from.

The gradient with respect to $z_i$ receives contributions from fish $i$'s rows
\emph{only}. If fish $i$ turns more often than the shared body predicts, this
gradient lowers $B_\theta z_i$, moving $z_i$ along $-B_\theta$. Each fish is
pushed independently of the others.

The gradient with respect to $B$ receives contributions from \emph{all} rows,
each weighted by its own fish's latent. A loading therefore grows only where
there is a consistent pattern across individuals --- where fish with large $z$
systematically require a larger forward logit. Had fish differed in unrelated
ways, the $z_i$ would point in unrelated directions and this gradient would
cancel. The axes reported above are in this sense discovered rather than
imposed.

\paragraph{Where the between-fish spread comes from.} At generation
$z\sim\mathcal{N}(0,I)$ and the applied offset is the linear map $Bz$. A linear
map of a standard normal has standard deviation equal to the norm of the map,
so
\begin{equation}
\mathrm{SD}\bigl(Bz\bigr) \;=\; \lVert B\rVert
\label{eq:sd-is-norm}
\end{equation}
exactly. The between-fish standard deviation of the model's population is
therefore just the \emph{length} of the loading vector, while its
\emph{direction} says which behaviours co-vary. For the travel channel,
$\lVert B_\ell\rVert=\sqrt{0.171^2+0.101^2}=0.199$, so individuals differ in
median travel distance by a factor $e^{0.199}=1.22$ per standard deviation.

\subsubsection{Why the prior does not set the scale}

Equation~\eqref{eq:sd-is-norm} makes the calibration question precise: what
fixes $\lVert B\rVert$? One might expect the prior to, since it penalises large
$z$. It does not, and the reason is structural.

Write the scalar analogue. Let $\delta_i$ be the offset that fish $i$'s data
requires, estimated from $n$ bouts with noise $\sigma^2$, and put
$\kappa=n/\sigma^2$:
\begin{equation}
\mathcal{L} \;=\; \frac{\kappa}{2}\sum_i\bigl(\delta_i-Bz_i\bigr)^2
              \;+\; \frac{1}{2}\sum_i z_i^2 .
\end{equation}
Setting both derivatives to zero,
\begin{align}
\frac{\partial\mathcal{L}}{\partial z_i}=0
  &\;\Longrightarrow\;
  z_i=\frac{\kappa B}{\kappa B^2+1}\,\delta_i=\frac{a\,\delta_i}{B},
  \qquad a=\frac{\kappa B^2}{\kappa B^2+1},
  \label{eq:zstar}\\[2pt]
\frac{\partial\mathcal{L}}{\partial B}=0
  &\;\Longrightarrow\;
  B=\frac{\sum_i z_i\delta_i}{\sum_i z_i^2}
   =\frac{(a/B)\sum_i\delta_i^2}{(a^2/B^2)\sum_i\delta_i^2}
   =\frac{B}{a}.
  \label{eq:Bstar}
\end{align}
Equation~\eqref{eq:Bstar} forces $a=1$, that is $\kappa B^2\to\infty$: at the
joint optimum the prior is inactive. The reason is visible in the loss --- only
the product $Bz$ is identified, and the penalty charges for a large $z$ but
nothing for a large $B$, so the optimiser escapes the prior by shrinking $z$
and growing $B$ in exact compensation. Substituting $a=1$ into
\eqref{eq:zstar} gives $Bz_i=\delta_i$: the fitted offset is the
\emph{unshrunk} per-fish estimate.

The consequence is a bias, and it runs in the opposite direction to the one
usually associated with random effects. Since $\delta_i$ is measured rather
than known, $\delta_i=\delta_i^{\text{true}}+\varepsilon_i$, and
\begin{equation}
\lVert B\rVert^2=\operatorname{Var}(\delta)
 =\underbrace{\operatorname{Var}(\delta^{\text{true}})}_{\text{wanted}}
 +\underbrace{\operatorname{Var}(\varepsilon)}_{\text{estimation noise}}
 \;>\;\operatorname{Var}(\delta^{\text{true}}) .
\end{equation}
This is the familiar upward bias of a variance component computed from
unshrunk per-group estimates: the fitted population is too \emph{wide}, not too
narrow. Consistent with this, the calibration below selects a gain below one,
contracting the loadings.

\subsubsection{Calibrating the scale}

A single global gain $g$ multiplies all three loadings. Because the induced
spread \eqref{eq:sd-is-norm} is exactly linear in $g$, choosing $g$ is
literally choosing the length of the loading vector:

\begin{center}
\begin{minipage}{0.86\textwidth}
\begin{verbatim}
for g in {0.4, 0.7, 1.2, 1.9, 3.0}:
    B  <- g * B_fitted
    generate a synthetic population of fish
    err(g) = sum over the 3 statistics of |log(SD_generated / SD_real)|
g* = argmin err(g)
\end{verbatim}
\end{minipage}
\end{center}

\noindent The three statistics are the across-fish spreads of mean $\log\ell$,
mean $\log\tau$ and forward fraction, computed on training fish only. The
objective is expressed in log-ratio units because the three live on different
scales, and an unnormalised sum would let the smallest of them dominate.

The gain cannot be chosen by likelihood. The training fish have their $z_i$
parked at $\delta_i/B$ by \eqref{eq:zstar}, so every one of them sits far out
along its own loading direction and a wider prior always scores them better;
the marginal likelihood is monotone in $g$ with no interior optimum. A
leak-free likelihood criterion would require refitting with each fish held out
in turn. This is the one population-level quantity not estimated by maximum
likelihood, and the model remains a normalised density throughout.

\paragraph{Re-centring.} Rescaling the loadings would move the pooled marginal
as a side effect, because for a log-normal channel a location shift enters the
arithmetic mean: $\E[e^{Bz}]=e^{\lVert B\rVert^2/2}$ at $z\sim\mathcal N(0,I)$.
Applying the gain therefore multiplies every generated travel and interval by
\begin{equation}
\exp\!\Bigl(\tfrac{1}{2}\bigl(g^2-1\bigr)\lVert B\rVert^2\Bigr)
\end{equation}
with nothing to cancel it. That exact quantity is subtracted from the
corresponding head biases when the gain is applied, so changing the width of
the population leaves its centre where it was. The turn loading needs no
analogue: it tilts a mixture logit rather than a location, so it is bounded by
the softmax and its pooled drift is held by the gain objective directly.

\subsubsection{Scoring an unseen fish}

A held-out fish has no fitted latent, so the model hedges across candidate
individuals on a Gauss--Hermite grid of $Q=49$ nodes, reporting at each bout
the predictive density averaged over its current belief and then updating that
belief:

\begin{center}
\begin{minipage}{0.86\textwidth}
\begin{verbatim}
log w[q] <- log prior(q)                      for q = 1..Q
for each bout t:
    L[t,q]  <- log p(y_t | x_t, z = node_q)   for q = 1..Q
    report     logsumexp(log w + L[t,.]) - logsumexp(log w)
    log w   <- log w + L[t,.]
\end{verbatim}
\end{minipage}
\end{center}

\noindent Read plainly: the model does not know which individual it is
watching, so it predicts by averaging over the individuals it still considers
possible, and re-weights them after every bout. Early bouts are predicted by a
diffuse average over the population; by the end of a trajectory the weight has
concentrated and predictions come from a nearly identified individual. That
progressive identification is where the likelihood improvement arises, and it
is unavailable to a model without a latent.

Because $z$ is constant along a trajectory, the belief update is a running
product and hence a cumulative sum in logs rather than a recursion, which is
why $49$ nodes cost little. The reported quantity is the same marginal
predictive density any latent-variable model reports, so Metric~A remains
comparable across models.

\paragraph{Generation.} A latent is drawn once per trajectory from the prior
and folded into the three head biases before the rollout begins; since it never
changes, the per-bout cost is zero. Drawing from the prior rather than from the
fitted table is what makes the generated fish a \emph{new} individual from the
estimated population rather than a reconstruction of the held-out one, and it
is the only choice that does not condition the generator on the data it is
about to be compared against.

% =====================================================================
